% SPDX-License-Identifier: CC-BY-4.0
\section{Limitations}
\label{sec:limitations}

\paragraph{Independent-Gaussian assumption.}
\sediment{} models each feature as an independent Gaussian under the
clean null. Multimodal feature distributions and feature-feature
correlations both fall outside this assumption. Mixture-aware
estimators and full-covariance robust estimators are natural
follow-ups. The \textsc{Anomaly-Within} ablation already shows a
rudimentary multivariate filter; a principled robust covariance
estimator is left to future work.

\paragraph{Single-host real-data run.}
The controlled poisoning evaluation runs on synthetic Gaussian
workloads. \cref{sec:eval-real-data} adds a single-host real-data
demonstration on captured Linux \texttt{bash} activity, but a
controlled multi-host evaluation across diverse workloads remains
future work. The synthetic setup is stronger than it looks: the
Gaussian assumption is the most favorable case for \textsc{Naive},
since its modeling assumption matches the data-generating
distribution. The gap we observe between \textsc{Naive} and
\sediment{} under poisoning is therefore a lower bound on what
would appear under real heavier-tailed workloads where \textsc{Naive}'s
outlier sensitivity is more pronounced.

\paragraph{Identity assumption.}
The threat model assumes the adversary cannot manipulate kernel-side
identity (PID, UID, cgroup). This holds for unprivileged adversaries
on properly-configured hosts but starts to fray under kernel-level
compromise. \chronosynd{} treats kernel-level adversaries as out of
scope.

\paragraph{Per-process scope.}
\sediment{} fits one baseline per process key, where the key in our
implementation is the process command name. An adversary who can
spawn a fresh process under an unfamiliar name will not have a
baseline for that process to drift against, only the ``never seen
this process before'' signal that the daemon surfaces as an
\texttt{[err]}-prefixed cache-miss line. Whether such cache misses
warrant an alert is an operator policy choice orthogonal to
\sediment{}'s correctness. In production it depends on the host's
expected process diversity. Our prototype emits the cache-miss
signal but takes no policy position on it.

\paragraph{Research artifact, not a production HIDS.}
\chronosynd{} is a research prototype. The algorithm, the BPF
runtime, the storage layer, the scoring engine, and the parity
checks are all real and tested, but service integration (systemd
units, alert delivery to syslog or SIEM, log rotation), baseline
lifecycle management (refit cadence, drift detection across the
weeks-to-months timescale of legitimate concept drift), performance
benchmarking under fleet-realistic load, CO-RE-based portable BPF
(one binary across heterogeneous kernel versions), per-PID rather
than per-\texttt{comm} baseline keying, and self-protection against
privileged tampering are all deliberately out of scope for this
paper. The artifact establishes that the algorithm is correct, the
runtime works on real Linux, and the cross-language parity is
defensible. Productionizing it for a SOC environment would take
several weeks of focused engineering and is outside the contribution
this paper claims.
